\section*{\sffamily\color{accent}Executive Summary}
\addcontentsline{toc}{section}{Executive Summary}

This report documents a Kolmogorov-Arnold Network Neural ODE (KAN-ODE)
\cite{koenig2024kanode,liu2024kan,chen2018neuralode} project that
reproduces the base paper's Lotka-Volterra benchmark, extends it to two
additional dynamical systems (a damped pendulum and an SIR epidemic model),
diagnoses and fixes two genuine cross-domain training failures, and runs
five independent follow-on studies probing solver order, differentiation
method, basis-function design, damping (originally intended as a stiffness
probe), and comparison against sparse symbolic regression.

\begin{keypoint}[Headline results]
\begin{itemize}
\item At the 10{,}000-epoch budget, KAN-ODE's extrapolation error is
  essentially flat out to $t=28$ ($1.1\times$ degradation), while a
  parameter-matched MLP degrades $8.3\times$ --- a $40.6\times$ KAN-ODE
  advantage at the longest horizon. \textbf{This advantage depends on the
  budget}: retrained to 50{,}000 epochs (with gradient clipping, on a
  different platform, Section~\ref{sec:phase4}), the same MLP extrapolates $6.8\times$ (near) to
  $10.8\times$ (far) \emph{better} than KAN-ODE.
\item The paper's claim that KAN-ODE needs $10\times$ fewer epochs than an
  MLP is not reproduced: the parameter-matched MLP reaches $10^{-3}$
  training loss $4.1\times$ sooner. The paper's literal $[2,50,2]$ tanh MLP
  fails to train under this recipe at any budget tested, and an audit of the
  paper's code identifies a learning-rate and initialization difference as
  the likely cause.
\item The pendulum failure (extrapolation $R^2=-4.64$ on the common test
  window, $-1.318$ on its own longer one) was fixed by extending the training
  window, to $R^2=+0.647$. SIR's divergence traced to a roughly $18$--$19\times$ unit-scale
  mismatch at initialization: a change of time units and a conservation
  projection remove the divergence ($R^2$ from $-20.97$ to $-0.145$), and a
  vanishing-field gate on the infected fraction, a physical prior, brings
  extrapolation $R^2$ to $+0.9981$.
\item Adjoint sensitivity already uses less memory than direct autograd at
  every trajectory length tested ($17.1$~MB against $18.6$--$21.5$~MB for
  the two training runs, and up to $\sim$160~MB at 3{,}201 points), but costs
  $1.9$--$2.1\times$ wall-clock --- direct autograd remains the correct
  default at this project's scale.
\item Gradient-norm roughness does not explain Euler's convergence floor
  (Spearman $\rho=-0.088$, $p=0.87$); shared, order-independent spike bursts
  dominate all six solvers' gradient traces.
\item A learnable spline/RBF basis blend never beats the better pure basis on
  Lotka-Volterra (it is $1.7$--$2.1\times$ worse there), or pure RBF on the
  pendulum, where no pure B-spline run was made. On
  the pendulum it overfits badly past $\sim$3{,}000 epochs, although measured
  directly the B-spline supplies only $\sim$1\% of its basis output.
\item The damping sweep never reaches a stiff regime (timescale ratio at
  most $4.3$), and no run diverged, as expected from the small step size
  ($|h\lambda|\le0.30$), which keeps every solver's discrete map close to the
  learned field's own flow; the only predicted instability is a mild
  numerical anti-damping for Euler at $\mu=0.1$. Its failures at
  $\mu\in\{0.5,1.0,2.0\}$ trace to the random seed rather than the solver,
  apart from one Euler-specific failure at $\mu=1.0$. At $\mu=2.0$, Newton's method
  confirms a genuine, attracting spurious equilibrium of the learned field,
  shared by all four solvers.
\item Sparse regression with finite-difference derivatives degrades far less
  under noise than KAN-ODE ($9\times$ vs.\ $1540\times$, $\sigma=0\to0.1$,
  from a much higher noise-free error) but is less accurate at every noise
  level. On noise-free
  data, given exact derivatives, it recovers the true Lotka-Volterra
  equation outright, so there the accuracy gap comes from the derivative
  estimate; whether it does under noise was not tested.
\item On a synthetic, non-mechanistic epidemic curve, two of SIR's three
  stability fixes transfer, including the one predicted not to.
\end{itemize}
\end{keypoint}

All seven basis functions, six integrators and the training loop were
implemented for this project and verified before use. Every integrator
reaches its theoretical order of convergence on a problem with a known
solution, and 249 of 250 automated tests pass (the remaining one, which needs a
GPU, is skipped). Every number and figure in this report is traceable to a saved
training record, a regenerable figure, or a post-hoc analysis script of a
saved checkpoint (Section~\ref{sec:verification}). The main limitations are
that almost every result is single-seed, that the budget reversal coincides
with a change in gradient clipping and in platform (operating system and
PyTorch version), and that three differences from the
paper's recipe (learning rate, initial-weight scale, RBF width) have been
identified but not yet measured.
